\documentclass[11pt]{article}
\usepackage{amsmath,amssymb,bm}
\usepackage{geometry}
\geometry{margin=1in}
\usepackage{setspace}
\usepackage{enumitem}
\usepackage{booktabs}
\usepackage{longtable}
\usepackage{hyperref}
\usepackage{xcolor}
\usepackage{listings}
\usepackage{titlesec}

\lstset{
  basicstyle=\ttfamily\small,
  breaklines=true,
  frame=none,
  columns=fullflexible
}

\titleformat{\paragraph}[runin]{\bfseries}{}{0em}{}[.]

\title{Progress Report:\\5-Layered LSIRM for Mixed-Type MIDUS Data}
\author{Hyunseok Yoon}
\date{April 2--6, 2026}

\begin{document}
\maketitle
\tableofcontents
\newpage

%======================================================================
\section{Overview}
%======================================================================

This report summarizes the development work carried out from April 2 to April 6, 2026, on the \textbf{5-layered nonhierarchical Latent Space Item Response Model (LSIRM)} for mixed-type data, applied to the MIDUS (Midlife in the United States) study.
The work encompasses four areas:
\begin{enumerate}[label=(\roman*)]
  \item Model specification refinements (layer-specific parameters, robust continuous layer),
  \item Ordinal layer reparameterization from cumulative link to Graded Response Model (GRM),
  \item Data pipeline expansion to include the MIDUS Refresher~1 cohort,
  \item Simulation studies comparing standard and robust formulations.
\end{enumerate}

%======================================================================
\section{Model Architecture}
%======================================================================

The model handles five data types simultaneously through a shared latent space.
All respondents share a single latent position $a_i \in \mathbb{R}^d$ across layers, while items have layer-specific positions $b_j^{(l)} \in \mathbb{R}^d$.

\begin{center}
\begin{tabular}{clll}
\toprule
Layer & Data type & Link & Parameters \\
\midrule
1 & Binary & Bernoulli--Logit & $\alpha_i^{(1)},\;\beta_j^{(1)},\;\gamma_1$ \\
2 & Continuous & Normal (robust: Student-$t$) & $\alpha_i^{(2)},\;\beta_j^{(2)},\;\gamma_2,\;\sigma_0^2,\;\lambda_{ij}^{(2)}$ \\
3 & Count & Negative Binomial & $\alpha_i^{(3)},\;\beta_j^{(3)},\;\gamma_3,\;\kappa$ \\
4 & Ordinal (5-pt) & GRM (v5) / Cumulative Link (v4) & $\alpha_i^{(4)},\;\gamma_4,\;\text{thresholds}$ \\
5 & Ordinal (4-pt) & GRM (v5) / Cumulative Link (v4) & $\alpha_i^{(5)},\;\gamma_5,\;\text{thresholds}$ \\
\bottomrule
\end{tabular}
\end{center}

The linear predictor for layers 1--3 follows:
\[
  \eta_{ij}^{(l)} = \alpha_i^{(l)} - \beta_j^{(l)} - \gamma_l \| a_i - b_j^{(l)} \|.
\]
For ordinal layers, there is no item difficulty $\beta_j$; the role of category-level difficulty is absorbed by the threshold parameters.


%======================================================================
\section{April 2: Layer-Specific Parameter Separation}
%======================================================================

\subsection{Layer-Specific Respondent Effects ($\alpha_i^{(l)}$)}

Previously, a single shared $\alpha_i$ was used across all five layers.
This was replaced by five independent respondent effects $\alpha_i^{(1)}, \dots, \alpha_i^{(5)}$, each updated via a separate Metropolis--Hastings block using only the corresponding layer's likelihood.

\paragraph{Rationale}
Different data types (binary symptoms, continuous biomarkers, Likert scales) have fundamentally different scales and respondent-level variation patterns.
A shared $\alpha_i$ forces all layers to compromise on a single respondent-level effect.

\paragraph{Implementation details}
\begin{itemize}[leftmargin=1.5em]
  \item Each $\alpha_i^{(l)}$ has its own proposal standard deviation: \lstinline{prop_sd$alpha1}, \dots, \lstinline{prop_sd$alpha5}.
  \item Each $\alpha_i^{(l)}$ has an independent prior variance $\sigma_{\alpha l}^2$.
\end{itemize}

\subsection{Layer-Specific Prior Variances ($\sigma_{\alpha l}^2$)}

The prior variance was similarly separated:
\[
  \alpha_i^{(l)} \sim N(0,\;\sigma_{\alpha l}^2), \quad
  \sigma_{\alpha l}^2 \sim \text{Inv-Gamma}(a_\sigma, b_\sigma), \quad l=1,\dots,5.
\]
Each $\sigma_{\alpha l}^2$ is updated via a conjugate Gibbs step:
\[
  \sigma_{\alpha l}^2 \mid \alpha^{(l)} \sim \text{Inv-Gamma}\!\left(a_\sigma + \frac{n}{2},\; b_\sigma + \frac{1}{2}\sum_{i=1}^n (\alpha_i^{(l)})^2\right).
\]

\paragraph{Files modified}
\begin{itemize}[leftmargin=1.5em]
  \item \texttt{my\_LSIRM\_5layered\_nonhierarchical\_v3.cpp}
  \item \texttt{my\_LSIRM\_5layered\_nonhierarchical\_cpp.R}
  \item \texttt{MIDUS\_5layered\_result.R}
\end{itemize}


%======================================================================
\section{April 5: Robust Continuous Layer (v4)}
%======================================================================

\subsection{Motivation}

Continuous biomarker variables (e.g., IL-6, CRP, cortisol) frequently exhibit heavy tails and outliers.
A Gaussian likelihood can be unduly influenced by extreme observations, biasing latent position estimates.

\subsection{Formulation}

We replace the Gaussian likelihood with a Student-$t$ distribution via a normal--gamma scale mixture:
\begin{align}
  \lambda_{ij}^{(2)} &\sim \text{Gamma}\!\left(\frac{\nu_2}{2},\;\frac{\nu_2}{2}\right), \\[4pt]
  Y_{ij}^{(2)} \mid \lambda_{ij}^{(2)} &\sim N\!\left(\mu_{ij}^{(2)},\;\frac{\sigma_0^2}{\lambda_{ij}^{(2)}}\right),
\end{align}
where $\mu_{ij}^{(2)} = \alpha_i^{(2)} - \beta_j^{(2)} - \gamma_2\|a_i - b_j^{(2)}\|$.

Marginalizing over $\lambda_{ij}^{(2)}$ yields $Y_{ij}^{(2)} \sim t_{\nu_2}(\mu_{ij}^{(2)}, \sigma_0^2)$.
The degrees of freedom $\nu_2$ is treated as fixed (default: $\nu_2=5$).

\subsection{MCMC Changes}

\begin{enumerate}[label=\textbf{Step \arabic*.}, leftmargin=3em]
  \setcounter{enumi}{-1}
  \item \textbf{(New)} Gibbs update for $\lambda_{ij}^{(2)}$:
  \[
    \lambda_{ij}^{(2)} \mid \cdots \sim \text{Gamma}\!\left(\frac{\nu_2+1}{2},\;\frac{\nu_2 + e_{ij}^2/\sigma_0^2}{2}\right),
    \quad e_{ij} = Y_{ij}^{(2)} - \mu_{ij}^{(2)}.
  \]
  \item All Layer-2 MH steps (for $\alpha_i^{(2)}, \beta_j^{(2)}, \gamma_2, a_i, b_j^{(2)}$) use the effective standard deviation $\sqrt{\sigma_0^2/\lambda_{ij}^{(2)}}$.
  \item The $\sigma_0^2$ Gibbs update uses weighted sum of squares: $\text{wSSE} = \sum_{i,j} \lambda_{ij}^{(2)} e_{ij}^2$.
\end{enumerate}

\subsection{Diagnostics}

Three diagnostic plots were added for the latent scale parameters:
\begin{itemize}[leftmargin=1.5em]
  \item Per-edge traceplot: 12 randomly selected $(i,j)$ pairs.
  \item Posterior mean heatmap: full $n \times P_2$ matrix; low $\lambda$ (red) indicates potential outliers.
  \item Per-item boxplot: distribution of $\lambda_{ij}^{(2)}$ by item, with $\lambda=1$ reference line.
\end{itemize}

\paragraph{Files created}
\begin{itemize}[leftmargin=1.5em]
  \item \texttt{my\_LSIRM\_5layered\_nonhierarchical\_v4.cpp}
  \item \texttt{my\_LSIRM\_5layered\_nonhierarchical\_cpp\_v4.R}
  \item \texttt{MIDUS\_5layered\_result\_v4.R}
\end{itemize}


%======================================================================
\section{April 5: Model Specification Document}
%======================================================================

A standalone LaTeX document (\texttt{model\_v2.tex}) was prepared, containing:
\begin{itemize}[leftmargin=1.5em]
  \item Full notation table.
  \item Likelihood specification for all five layers.
  \item Complete prior hierarchy.
  \item 9-step MCMC algorithm using \texttt{algorithm2e} environment, with MH vs.\ Gibbs steps clearly distinguished.
\end{itemize}


%======================================================================
\section{April 5: 4-Layered Robust Model \& Simulation}
%======================================================================

\subsection{4-Layered Robust Model (v5)}

The robust continuous layer was back-ported to the 4-layered model variant (which uses a single shared $\alpha_i$ and one ordinal layer) for simulation studies.

\begin{center}
\begin{tabular}{lcc}
\toprule
Feature & 4-layered (v5) & 5-layered (v4) \\
\midrule
Alpha & Shared $\alpha_i$ & Independent $\alpha_i^{(1)},\dots,\alpha_i^{(5)}$ \\
$\sigma_\alpha^2$ & Single & 5 independent \\
Ordinal layers & 1 & 2 (5-pt, 4-pt) \\
Robust L2 & \checkmark & \checkmark \\
\bottomrule
\end{tabular}
\end{center}

\subsection{Simulation Study: Standard vs.\ Robust}

A simulation study compares the standard (Gaussian L2) and robust (Student-$t$ L2) models.

\paragraph{Design}
\begin{itemize}[leftmargin=1.5em]
  \item \textbf{Scenario 1} (Standard): Data generated with Student-$t$ errors ($\nu=5$).
  \item \textbf{Scenario 2} (Outlier): Additionally, 10\% of respondents have latent positions drawn from $N(0,8)$ (extreme outliers).
  \item 100 replications, parallel execution via \texttt{mclapply}.
  \item MCMC: 50{,}000 iterations, 10{,}000 burn-in, thin by 10.
\end{itemize}

\paragraph{Evaluation}
95\% credible interval coverage for: $\alpha$, $\beta_1$--$\beta_3$, $u$, $\delta$, and layer-specific distances $d_{ij}^{(l)} = \gamma_l \|a_i - b_j^{(l)}\|$.


%======================================================================
\section{April 6: MIDUS Refresher~1 Data Integration}
%======================================================================

\subsection{Background}

The MIDUS 2 (2004--2006) sample provides limited sample sizes after screening for depression criteria ($\text{screen\_case}=1$) and requiring complete cases across biomarker and cognitive sub-projects (e.g., 148 for the full P1--P3--P4 analysis).
To increase statistical power, we incorporated the MIDUS Refresher~1 cohort (2011--2014), which independently recruited 3{,}577 new participants.

\paragraph{Why Refresher, not Wave~3?}
MIDUS~3 (2013--2014) is a pure longitudinal follow-up of MIDUS~2 participants.
Adding Wave~3 data for individuals already in Wave~2 would introduce within-subject dependence.
Since MIDUS~3 contains no new recruits, the only non-overlapping individuals would come from sub-project participation differences---introducing selection bias.
The Refresher cohort, being independently recruited, avoids both issues.

\subsection{Variable Mapping}

The Refresher uses different variable prefixes but measures the same constructs:

\begin{center}
\begin{tabular}{lll}
\toprule
Category & Wave~2 prefix & Refresher~1 prefix \\
\midrule
Survey & \texttt{B1PA*}, \texttt{B1SA*} & \texttt{RA1PA*}, \texttt{RA1SA*} \\
Biomarker & \texttt{B4*} & \texttt{RA4*} \\
Cognitive & \texttt{B3*} & \texttt{RA3*} \\
ID variable & \texttt{M2ID} & \texttt{MRID} \\
\bottomrule
\end{tabular}
\end{center}

One variable (\texttt{B4BSCL4A}, cortisol final average) was absent in the Refresher data.
It was removed from the Wave~2 pipeline as well to ensure structural alignment (7 continuous variables retained).

\subsection{Data Combination}

Wave~2 and Refresher~1 are combined by row-binding the response matrices ($Y_\text{bin}$, $Y_\text{con}$, $Y_\text{cnt}$, $Y_\text{ord1}$, $Y_\text{ord2}$), with column names standardized to Wave~2 conventions.
A \texttt{source} indicator (\texttt{"wave2"} / \texttt{"refresher1"}) is maintained for potential use as a covariate.

\paragraph{Files created/modified}
\begin{itemize}[leftmargin=1.5em]
  \item \texttt{MIDUS\_preprocess\_refresher\_v3.R} (new)
  \item \texttt{MIDUS\_preprocess\_2\_v3.R} (modified: \texttt{B4BSCL4A} removed)
  \item \texttt{MIDUS\_5layered\_result\_v4.R} (modified: combined data pipeline)
\end{itemize}


%======================================================================
\section{April 6: GRM Ordinal Layer (v5)}
%======================================================================

\subsection{Motivation}

The original ordinal layers (L4, L5) used a cumulative link model with a softplus reparameterization to enforce threshold ordering:
\[
  c_{j,1} = 0, \quad c_{j,k+1} = c_{j,k} + \text{softplus}(\delta_{j,k}), \quad k=1,\dots,K-2.
\]
Thresholds were updated jointly (all $\delta_{j,\cdot}$ simultaneously) via MH with a Jacobian correction.

This was replaced by the \textbf{Graded Response Model (GRM)} parameterization, following the implementation in the \texttt{lsirm12pl} R~package.

\subsection{GRM Formulation}

The GRM defines cumulative probabilities in the \emph{exceedance} direction:
\[
  P(Y_{ij} \ge k \mid \eta_{ij}) = \text{logistic}(\eta_{ij} + \beta_{j,k}), \quad k=1,\dots,K-1,
\]
where $\eta_{ij} = \alpha_i^{(l)} - \gamma_l \|a_i - b_j^{(l)}\|$ and the thresholds satisfy $\beta_{j,1} > \beta_{j,2} > \cdots > \beta_{j,K-1}$ (descending order).

The category probabilities are:
\[
  P(Y_{ij}=k) = \begin{cases}
    1 - P(Y \ge 1) & k=0, \\
    P(Y \ge k) - P(Y \ge k+1) & 1 \le k \le K-2, \\
    P(Y \ge K-1) & k=K-1.
  \end{cases}
\]

\subsection{Key Differences from v4}

\begin{center}
\begin{tabular}{lll}
\toprule
Aspect & v4 (Cumulative Link) & v5 (GRM) \\
\midrule
Threshold structure & Ascending, softplus-linked & Descending, direct \\
Parameters & $\delta_{j,k}$ (unconstrained) & $\beta_{j,k}$ (order-constrained) \\
Number of params & $P \times (K-2)$ & $P \times (K-1)$ \\
Update strategy & Joint MH + Jacobian & Individual MH + constraint check \\
Data encoding & 1-based internally & 0-based internally \\
\bottomrule
\end{tabular}
\end{center}

\subsection{Threshold Update Algorithm}

For each item $j$ and each threshold index $c = 1,\dots,K-1$:
\begin{enumerate}[leftmargin=2em]
  \item Propose $\beta_{j,c}^* \sim N(\beta_{j,c}, \sigma_\text{prop}^2)$.
  \item \textbf{Check ordering constraint}: reject immediately if $\beta_{j,c-1} \le \beta_{j,c}^*$ (when $c>1$) or $\beta_{j,c}^* \le \beta_{j,c+1}$ (when $c<K-2$).
  \item Compute the MH ratio using the GRM log-likelihood summed over all respondents and a $N(\mu_\beta, \sigma_\beta^2)$ prior.
  \item Accept or reject.
\end{enumerate}

This mirrors the approach in \texttt{lsirm12pl::lsirm\_grm.cpp}, adapted to our multi-layer shared latent space framework.

\paragraph{Files created}
\begin{itemize}[leftmargin=1.5em]
  \item \texttt{my\_LSIRM\_5layered\_nonhierarchical\_v5.cpp}
  \item \texttt{my\_LSIRM\_5layered\_nonhierarchical\_cpp\_v5.R}
  \item \texttt{MIDUS\_5layered\_result\_v5.R}
\end{itemize}


%======================================================================
\section{Summary of Code Versions}
%======================================================================

\begin{center}
\begin{tabular}{clp{8cm}}
\toprule
Version & Date & Description \\
\midrule
v3 & Apr 2 & Layer-specific $\alpha_i^{(l)}$, $\sigma_{\alpha l}^2$; cumulative link ordinal \\
v4 & Apr 5 & + Robust continuous layer (Student-$t$ via Normal--Gamma mixture) \\
v5 & Apr 6 & + GRM ordinal parameterization (replacing cumulative link) \\
\bottomrule
\end{tabular}
\end{center}

Each version consists of three files:
\begin{enumerate}[leftmargin=1.5em]
  \item C++ MCMC engine: \texttt{my\_LSIRM\_5layered\_nonhierarchical\_v\{N\}.cpp}
  \item R wrapper (initialization + Procrustes): \texttt{my\_LSIRM\_5layered\_nonhierarchical\_cpp\_v\{N\}.R}
  \item Analysis runner (data loading + MCMC + diagnostics): \texttt{MIDUS\_5layered\_result\_v\{N\}.R}
\end{enumerate}


%======================================================================
\section{Current Model Summary (v5)}
%======================================================================

The current model (v5) combines all developments:
\begin{itemize}[leftmargin=1.5em]
  \item \textbf{5 layers}: binary (L1), continuous--robust (L2), count (L3), ordinal-5pt (L4), ordinal-4pt (L5).
  \item \textbf{Shared latent positions} $a_i$ across all layers; layer-specific item positions $b_j^{(l)}$.
  \item \textbf{Layer-specific respondent effects} $\alpha_i^{(l)}$ with independent prior variances $\sigma_{\alpha l}^2$.
  \item \textbf{Layer-specific distance weights} $\gamma_l$.
  \item \textbf{Robust continuous layer}: Student-$t$ errors via Normal--Gamma mixture with $\lambda_{ij}^{(2)}$ latent scales.
  \item \textbf{GRM ordinal layers}: direct descending thresholds with individual constrained MH updates.
  \item \textbf{Data}: MIDUS Wave~2 + Refresher~1 combined (independent cohorts, no ID overlap).
\end{itemize}

\end{document}
